\section{Background and Related Work}
\label{sec:ch5-background-related}
We situate this work at the intersection of biophilic design, interaction design (IxD), and research on indoor environments. We review foundational perspectives on biophilic design in architecture. We then examine how biophilic principles have been translated into indoor environments within HCI. Finally, we examine how technologically-mediated indoor environments engage with experiential concerns and position our contribution within speculative, design-led HCI.

\subsection{Biophilic Design in Architecture}

Biophilic design in architecture is grounded in the proposition that humans have an inherent orientation toward living systems and lifelike processes~\cite{wilson1986biophilia}. This orientation extends beyond direct contact with nature to encompass emotional, cognitive, and ecological forms of awareness, influencing wellbeing, attention, and environmental meaning-making~\cite{ulrich1984view, beatley2009biophilic}. Architectural interpretations of biophilia have therefore focused on how spatial organisation, material choices, and environmental changes can sustain ongoing connections between occupants and natural processes~\cite{gil2015biophilic, aristizabal2021biophilic}. Historically, architectural traditions have expressed biophilic sensibilities through environmental responsiveness, and attention to change over time, for example in climate-adaptive timber architecture and ecological planning, that treat buildings as part of larger environmental systems~\cite{beatley2016chapter, kim2009alvar, mcharg1969design}.

These sensibilities were later formalised into biophilic design principles. \citet{kellert2018nature} articulated attributes such as access to daylight, natural materials, prospect and refuge~\cite{dosen2016evidence}, and indirect references to nature as strategies integrated into architectural practice~\cite{kellert1995biophilia, kellert2018nature, kellert2008dimensions}. More recently, biophilic architecture has extended this trajectory through climate-adaptive envelopes, passive ventilation, daylight-driven spatial organisation, and landscape integration, motivated by sustainability and long-term environmental resilience~\cite{tekin2023impact}. Recent work in architecture-adjacent domains, particularly building science, has developed human-centred frameworks examining how occupant experience, behaviour, and evaluation relate to spatial configuration, environmental conditions, and technological integration in indoor environments~\cite{khan2025building, roofigari2025conceptual, economidou2024lived}. These approaches emphasise wellbeing, personalised environments, and adaptive systems (e.g. HVAC, lighting, IoT) that align building performance with human needs at scale. However, they largely retain an optimisation-oriented logic, framing experience in terms of measurable and adjustable parameters, rather than ongoing sensory engagement or environmental variation~\cite{becerik2022ten}.
% These approaches position biophilic design as an architectural commitment to embedding ecological sensibilities and sensory variation into everyday built environments.


\subsection{Biophilic Design in HCI and IxD}
In HCI, biophilic design has inspired a range of interactive systems aiming to support connection with nature, restoration, or environmental awareness. These include sunlight responsive lighting~\cite{wang2015intelligent}, soft robotics that represent nature~\cite{chooi2022symbiosis, sabinson2021plant, sabinson2024every}, or projected nature scenes~\cite{sabinson2023walk}. Some use metaphorical interfaces to encourage mindfulness~\cite{wang2022nature, jung2024leaf}; others embed tangible cues into everyday routines~\cite{nam2023florawear, loh2010please}. In both cases, interaction is commonly framed as episodic engagement with a designed biophilic interface. While most systems engage visual and auditory modalities, tactile, olfactory, and proprioceptive channels remain under-explored~\cite{dias2024quantifying, yildirimmultisensory}. Multisensory layering is rarely treated as a design material in its own right. Biophilic interactivity is often framed around individual wellbeing, playfulness, or affect regulation, with limited engagement with relational, collective, or situated modes of inhabitation.



\subsection{Interactions in Built Environments}
A substantial body of work focuses on comfort as a primary experiential objective~\cite{alavi2017comfort, alavi2020five}. This includes personalised thermal comfort systems~\cite{jazizadeh2014human, khare2023thermal}, interactive demonstrations of airflow~\cite{zhong2020hilo, snow2019performance, gaver2013indoor}, as well as acoustic~\cite{marques2020real, choi2009decision} and visual comfort optimisation~\cite{bian2023simulation, zheng2023daylighting}. Occupants are typically the recipients or regulators of system behaviour, with few works looking at social negotiation and shared interactions~\cite{clear2018thermokiosk, von2021towards, clear2017d}.

Beyond comfort, work explores how interactive technologies shape experiential, affective, and social engagement with indoor environments. Ambient displays and media-based interventions use light, sound, or projection to evoke calming or emotionally resonant atmospheres~\cite{feng2019livenature, behrens2018sentiment, rao2023towards}.  Discovery is also explored through installations that invite playful, interpretive engagement rather than task completion~\cite{gaver2006history, wiberg2020interaction}. Examples such as \emph{Pendaphonics} foreground bodily movement and collective improvisation through spatial interaction~\cite{hansen2009pendaphonics}. Work on proximity and spatial movement highlights how interaction emerges through bodily orientation and movement, such as the ``wifi-suit'' that reacts to the strength of wireless signals as the wearer moves around the office~\cite{van2018skin}. Others examine neurodiverse and sensory-sensitive design that accounts for varied perceptual needs~\cite{swaminathan2021tactile, garzotto2020interactive}. We aim to broaden this experiential scope of interactive indoor environments by introducing biophilic interactive experiences.  

\subsection{Speculative and Design-Led Approaches for Interactive Experiences in Built Environments}
This work adopts a speculative, design-led methodology. In HCI, speculative approaches have been used to explore future interaction trajectories by situating technologies within imagined contexts, through prototypes, scenarios, or narratives 
that retain ambiguity to provoke interpretation and critique~\cite{galloway2018speculative, lee2025speculative}.

Within built environment research, prior work has also  employed sketches and scenarios to reason about interaction in energy and building management systems~\cite{rodden2013home, pschetz2019autonomous}. HBI has further used participatory and speculative formats to surface assumptions about smart environments, including student-led design explorations that revealed concerns around automation and loss of control~\cite{mitchell2020no}, and public-facing design fiction events that foregrounded everyday expectations, scepticism, and values around smart building futures~\cite{paraschivoiu2023postcards, economidou2023engaging}. Related exploratory studies of quantified buildings similarly examine how occupants interpret sensing and automation in practice, highlighting tensions between optimisation and lived experience~\cite{margariti2023understanding}. We adopt an expert-led design session to capture expert imaginaries and thinking on future biophilic interactive experiences in built environments. 

\subsection{Gap and Positioning}
Existing biophilic IxD work focuses on visual metaphors, restorative effects, and individual wellbeing outcomes, often treating interaction as a means to deliver calming or nature-inspired experiences. While valuable, this emphasis risks overlooking other ways of relating to nature indoors. This paper builds on architectural perspectives while complementing existing occupant-centred frameworks by articulating biophilic interaction at the level of interaction design. We foreground how biophilic qualities can be composed, valued, and enacted through interaction at architectural scale. By developing a \textbf{biophilic interaction design framework} comprising dimensions, values, and interaction forms grounded in expert reasoning, we broaden interaction design for indoor environments and provide a foundation for future design-led and empirical research.
